\documentclass[12pt,a4paper,twoside]{scrreprt}

% Sprachpakete
\usepackage[utf8]{inputenc}
%\usepackage[latin1]{inputenc}
\usepackage[T1]{fontenc}
\usepackage{lmodern}
\usepackage[english]{babel}
\usepackage{url}
\usepackage{hyperref}
% Zum plattformunabhängigen Einfügen anderer Tex-Dateien
% Hinweise zur Benutzung weiter unten
\usepackage{subfiles}
% Zum Anpassen der Seitenränder
\usepackage{geometry}
% Einrücken vermeiden
\usepackage{parskip}
% Zum Einbinden von PDF-Seiten
\usepackage{pdfpages}
% Zum Einbinden von Grafiken
\usepackage{graphicx}		
\usepackage{morefloats}
% Um Zellen in Tabellen zusammenzufassen
\usepackage{multirow}
% Für farbige Schriften und Tabellen
\usepackage{xcolor}
\usepackage{colortbl}
% Package für zusätzliche mathematische Symbole
\usepackage{amssymb}
\usepackage{units} 
% Package zum erstellen von ToDo-Anmerkungen
% Befehl: \todo{Text} oder \todo[inline]{text} für Anmerkungen am Rand oder in-line
%					\missingfigure{text} Platzhalter für geplante Grafiken/Tabellen mit Beschreibung
% Beim fertigen Dokument muss dieser Package mittels \usepackage[disable]{todonotes} deaktiviert werden,
% um alle vorhandenen ToDos als auch die ToDo-Liste ausgeblendet werden.
\usepackage[disable]{todonotes}	

% Nummerierung der Grafiken nach Kapitel
\usepackage{chngcntr}
\counterwithin{figure}{section}


%%%%%%%%%%%%%%%%%%%%%%%%%%
\begin{document}

\pagestyle{plain}

%%%%%%%%%%%%%%%%%%%%%%%%%%%%%%%%%%%%%%%%%%%%%%%%%%%%%%%%%%%%%%%%%%%%%%%%%%%% Titleseite V1
\begin{titlepage}

		\begin{center}
		\includegraphics[scale=0.3]{LogoUniOL.png}\\[3cm]
		\Huge{\textbf{Projektgruppe Bull's Eye}}\\[0.4cm]
		\huge{\textbf{Cardboard assembling}}\\[0.5cm]
		\Large{2016~-~2017}\\[3ex]
		\large{
		Daniela Betzl\\
		Tim Lukas Cofala\\
		Aljoscha Niazi-Shahabi\\
		Stefan Niewerth\\
		Henrik Reichmann\\
		Rieke von Bargen\\
		}
		\vspace{1.5cm}
		\large{
				Prof. Dr. Susanne Boll-Westermann\\
				Advisors: M.Sc. Tim Stratmann, M.Sc. Uwe Gr\"unefeld\\
				Organization: M.Sc. Marion Koelle \\
				Last revision: 26.09.2017\\
				Version: 2.0\\
				Carl von Ossietzky University of Oldenburg\\
				\textsc{Faculty II } School of Computing Science, Business Administration, Economics, and Law 
			}
		\end{center}
\end{titlepage}

% Diese Liste listet alle eingefügten ToDo-Anmerkungen auf, sodass man ein Überblick über allen ToDos im ganzen Dokument hat
\listoftodos

%%%%%%%%%%%%%%%%%%%%%%%%%%%%%%%%%%%%%%% Abstract
% Auskommentieren, falls kein Abstrakt notwendig ist
%\subfile{chapter/abtract.tex}

%%%%%%%%%%%%%%%%%%%%%%%%%%%%%%%%%%%%%%% Verzeichnisse
\pagenumbering{roman}
\tableofcontents \clearpage
%\listoffigures \clearpage
%\listoftables \clearpage
\pagenumbering{arabic}
\clearpage
\setcounter{page}{1}

%%%%%%%%%%%%%%%%%%%%%%%%%%%%%%%%%%%%%%% die Kapitel des Dokuments
% Die einzelne Kapitel können hier mittels dem Befehl \subfile{chapter/unterordner/textdatei.tex} eingefügt werden,
% wobei chapter/unterordner/ der relative Pfad zur Datei ist.
% Das Package subfiles sorgt dafür, dass die einzlenen Kapitel automatisch mit dieser Datei verknüpft werden, sodass keine 
% Konfiguration des LaTeX-Editors notwendig ist. Dies ist praktisch, da die Konfiguration im Editor manuell erledigt werden muss,
% und von Editor zu Editor sich unterscheiden.
% Zudem ermöglicht subfiles die einzelnen Kapitel direkt zu kompilieren, ohne masterdoc.tex sowie die anderen zugehörigen Tex-Dateien
% kompilieren zu müssen. Dadurch geht das kompilieren schneller, wobei nachdem man fertig ist masterdoc ebenfalls kompilieren sollte, 
% sodass das Kapitel bzw. die Änderungen am Kapitel im Hauptdokument ebenfalls zu sehen ist.

%%%%%%%%%%%%%%%%%%%%%%%%%%%% Kapiteln 
%\subfile{chapter/platine/platine.tex}
%\subfile{chapter/cutter/cutter.tex}
\subfile{chapter/template/template.tex}

\end{document}